% !TEX root = ../main.tex
%
% Methods subsection -- Data Preprocessing
% Requires (in preamble of main.tex): booktabs

\subsection{Data Preprocessing}
\label{subsec:preprocessing}

\subsubsection{Whole-Slide Image Tiling}

Whole-slide images were tiled at a target resolution of 1.0 micron per pixel
(MPP) into $1024 \times 1024$ pixel patches. Native slide MPP was read from
scanner metadata and validated to fall within a plausible range (0.05--5.0
$\mu$m/pixel) to catch corrupted or mislabeled metadata before tiling.
Background regions were excluded using a tissue mask computed by Otsu
thresholding the saturation channel of a low-resolution slide thumbnail; a
candidate tile was kept only if at least 15\% of its area fell within the
tissue mask. Retained tiles were passed to the frozen patch encoder at their
native $1024\times1024$ resolution, normalized with standard ImageNet channel
statistics.

\subsubsection{Clinical Data Harmonization}

TCGA and CPTAC clinical variables were mapped to a common schema. Overall
survival time was defined from diagnosis to death (event = 1) or to last
follow-up (censored, event = 0). Pathologic stage, TNM classification, and
tumor grade were deliberately excluded from all model inputs, since these are
summary variables already strongly associated with prognosis and would act as
shortcut features that bypass the imaging and transcriptomic signal being
evaluated; only age and sex are used as the clinical input.

\subsubsection{RNA-seq Processing}

TCGA-PAAD expression values (STAR-Counts) were transformed as
$\log_2(\text{FPKM-UQ} + 1)$; CPTAC-PDAC expression values were used as
provided (RSEM upper-quartile normalized, already on a $\log_2$ scale). Genes
were restricted to those common to both cohorts' processed matrices, and each
gene was $z$-score normalized independently within each cohort.

\subsubsection{Literature-Guided Gene Selection}

The full common-gene set is far larger than what a cohort of this size can
support without overfitting, and a subtype-classification gene panel is not
necessarily the panel most informative for survival. A compact,
survival-oriented gene set was therefore constructed. First, 163 genes were
curated from the pancreatic ductal adenocarcinoma (PDAC) literature across
eight functional categories (Table~\ref{tab:gene-categories}).

Second, for each cross-institution training direction (e.g.,
TCGA-PAAD $\rightarrow$ CPTAC-PDAC), a univariate Cox partial-likelihood score
test was computed for every common gene using only that direction's own
training-split cases (91 TCGA-PAAD patients, 50 observed deaths, for the
TCGA-PAAD $\rightarrow$ CPTAC-PDAC direction). The cohort that serves as the
external test set for that direction is not accessed at any point during gene
selection. This is a deliberate departure from combining both cohorts' training
splits (e.g., by Stouffer's method) into a single ranking: because a
substantial fraction of either cohort's patients fall within the training
split used for such a combined ranking, reusing a combined ranking for external
evaluation would leak survival information about patients later scored as
external test cases.

Third, the final gene ordering places the curated literature genes first
(ordered by their own within-direction Cox rank), with the remaining ordering
determined by genome-wide Cox score. Two criteria were used to cut this
ordering into a final gene set, and both are evaluated in Results: (i) a fixed
top-$N$ ($N \in \{500, 1500\}$), and (ii) a statistically derived cutoff based
on the Benjamini--Hochberg false discovery rate (FDR), retaining the curated
genes plus any additional gene with $q < 0.1$. Because approximately 20,000
genes are tested, an uncorrected $p<0.05$ cutoff would retain roughly 1,000
genes by chance alone; FDR correction is used instead of a raw $p$-value
threshold for this reason.

\begin{table}[htbp]
  \centering
  \begin{tabular}{@{}l c@{}}
    \toprule
    \textbf{Category} & \textbf{Genes} \\
    \midrule
    Core driver / tumor suppressor            & 16 \\
    DNA damage repair / therapy response       & 13 \\
    Classical / pancreatic progenitor          & 22 \\
    Basal-like / squamous / mesenchymal        & 21 \\
    Stroma / ECM / invasion                    & 30 \\
    Immune / inflammation / TGF-$\beta$        & 25 \\
    Proliferation / cell cycle / apoptosis     & 18 \\
    Hypoxia / metabolism / acinar program      & 18 \\
    \midrule
    Total curated genes                        & 163 \\
    \bottomrule
  \end{tabular}
  \caption{Literature-curated PDAC gene categories used to seed the
    survival-guided gene ranking. Individual gene symbols are listed in the
    Supplementary Information.}
  \label{tab:gene-categories}
\end{table}
